\chapter{Conclusões e Recomendações}
\label{cap:05}

% O NITE se encontra em fase de desenvolvimento, com avanços significativos de suas principais funcionalidades. Apesar do progresso, ainda se faz necessário dar continuidade ao projeto, reservando espaço para futuras implementações que irão contribuir para o funcionamento total do editor CLI. Ajustes referentes a interface e usabilidade devem seguir em todos os passos (incluindo os de estágio avançado), seguindo de levantamentos e apresentação de resultados, visando usabilidade. No decorrer dos proximos meses, pretende-se concluir sua fase de testes e desenvolvimento, possibilitando a apresentação dos resultados finais, impactos e propostas de soluções.

Ao longo de diversos estudos sobre ferramentas correlatas ao NITE e da biblioteca \textit{ncurses} para a linguagem C, do desenvolvimento da ferramenta e dos desafios encontrados, pode-se concluir que os objetivos propostos foram atingidos: o editor é funcional, minimalista e oferece ao usuário final modos distintos de interação (edição e leitura), priorizando aquilo que é essencial em uma interface de linha de comando. Ainda assim, diante da diversidade de ferramentas existentes e do potencial da aplicação, o NITE entra agora em uma nova etapa, aberta a futuras atualizações, melhorias de qualidade de vida e novas funcionalidades.